\documentclass[10pt,a4paper]{amsart}
\usepackage[margin=19mm]{geometry}
\usepackage{amsmath,amssymb,mathtools}
\usepackage[scr=rsfs,cal=euler]{mathalfa}
\usepackage{xfrac}
\usepackage[unicode,hidelinks,bookmarks=false]{hyperref}
\newcommand{\ie}{i.e.\ }
\newcommand{\cf}{cf.\ }
\newcommand{\resp}{respectively\ }
\newcommand{\Subsection}{\subsection}
\providecommand{\nobreakdash}{\nobreak\hbox{-}}
\setlength{\emergencystretch}{2em}
\multlinegap0pt
\def\CC{\mathbb{C}}
\usepackage{amssymb}
\newtheorem{proposition}{Proposition}
\title{Strong Ramadanov Conjecture for Real Ellipsoids in $\CC^2$: two approaches}
\author{Jan Gregorovic, Ilya Kossovskiy, Son-Ying Li,, Kevin Shi}
\date{Source: arXiv:2607.23185v1 (CC BY 4.0)}
\begin{document}
\maketitle

\noindent\emph{Source and licence.} Strong Ramadanov Conjecture for Real Ellipsoids in $\CC^2$: two approaches. Version arXiv:2607.23185v1 (CC BY 4.0), distributed by arXiv under the Creative Commons Attribution 4.0 International licence, http://creativecommons.org/licenses/by/4.0/, as recorded in the arXiv licence metadata for this exact version and shown on the abstract page https://arxiv.org/abs/2607.23185v1. Full licence notice: \texttt{licenses/arxiv-2607.23185-CC-BY-4.0.txt}.\par\smallskip

\noindent\textit{Editorial scope.} Counted unit: the article's proposition weightcheck (source0.tex lines 1526--1528). The article's complete original proof of it is delivered with it (source0.tex lines 1529--1540, one \texttt{proof} environment of 12 lines). No mathematics is written, repaired or re-derived: the statement and the proof are the article's own lines, in the article's own notation, and no retained line is substituted or re-flowed.  No further source-local context is needed: every cross-reference of every retained line resolves inside the two delivered spans.  Nothing was omitted from any retained span, so the package carries no omission record.  No retained line contains a citation, so the wrapper carries no bibliography and no bibliography entry.  No retained line contains a figure environment, an image-inclusion command or drawing code, so no image is delivered and the package declares no asset dependency.  Editorial formatting only, in addition: an AMS \texttt{amsart} preamble that restates the article's own declarations and macros used by the retained lines, namely the environments \texttt{proposition} on separate counters, together with the standard packages those lines need.  The retained environments print in this document as Proposition 1 in order of appearance, whereas the article prints the same items with the numbering of its own full document.  The complete native source of the article as distributed in the arXiv e-print for this exact version is bundled unchanged as \texttt{sources/206009/source0.tex}.
\begin{proposition}\label{weightcheck}
The Tresse invariant of weight $(r,s)$ is relative CR invariant of CR weight $(r+s,s)$. Moreover, the Tresse invariant of weight $(r,s)$ of the lowest order $k$ vanishes if and only if $N^{2s-k+r}_{1-s+k,1-s+k-r}$ vanishes.
\end{proposition}

\begin{proof}
Note that the CR weights define the weighting system:
$$
z\mapsto \lambda z,\quad w\mapsto \lambda \mu w,\quad p\mapsto \mu p
$$
(recall that $p=dw/dz$).
Thus $ \tilde  \lambda=\lambda$ and $\tilde \mu= \lambda \mu$ and the Tresse invariant of weigth $(r,s)$ is relative CR invariant of CR weight $(r+s,s)$.

Let us take the order of the invariant into account: $N^c_{a,b}$ is of order $a+b+c$, while $\Phi$ is obtained by two derivations, i.e., $a+b+c-2=k$. Among the Tresse invariant of weigth $(r,s)$ there is unique one of the lowest order $k$ and thus it is  $N^c_{a,b}$ with $a+b+c-2=k,a+c-1=r+s,b+c-1=s$ 
%[a = 1-s+k, b = 1-s+k-r, c = 2*s-k+r]
up to nowhere vanishing absolute invariant, i.e., it vanishes if and only if the Tresse invariant vanishes.
\end{proof}

\end{document}
